\chapter{Dynamic Linking, Loading, TLS, and Link-Time Transformations}
\label{ch:dynamic-linking}

A statically linked image can resolve every ordinary reference before execution.
A shared library may not know either its load address or the particular definition that will satisfy a globally visible name.
Its instructions must therefore cooperate with a runtime linker rather than assuming that all fields are final.
This changes symbol lookup, address formation, initialization, and ownership of the tables that contain writable addresses.
It also changes which optimizations a static linker can justify.

The Lua-mapped linker developed in this sequence emits static executables.
This chapter studies the additional runtime protocol using the ELF and AMD64 specifications and a separate executable experiment with the host dynamic loader~\cite{xinuos-2026-elf,matz-2012-amd64}.
Keeping that distinction explicit prevents a static relocation patcher from being mistaken for a shared-library implementation.

\section{Load bias and relative addressing solve different problems}
\label{sec:ch59-bias}

Suppose a shared object's linked virtual layout is loaded with bias $B$.
A symbol at linked virtual address $S_0$ then has runtime address $B+S_0$ when it binds within that object.
A PC-relative reference from $P_0$ has difference
\[
 (B+S_0)-(B+P_0)=S_0-P_0.
\]
The common bias cancels, so an internal relative reference can remain unchanged.
An absolute pointer stored in data still needs the bias incorporated somewhere.
Figure~\ref{fig:ch59-bias} distinguishes the unchanged displacement from the changed pointer.

\begin{figure}[htbp]
\centering
\begin{tikzpicture}[node distance=15mm and 23mm]
\node[basic block] (p) {place P}; \node[basic block,right=of p] (s) {symbol S};
\node[basic block,below=of p] (bp) {place B+P}; \node[basic block,right=of bp] (bs) {symbol B+S};
\draw[cfg edge] (p)--node[above=3mm,font=\small] {difference S-P}(s);
\draw[cfg edge] (bp)--node[below=3mm,font=\small] {same difference}(bs);
\draw[dashed,->] (p)--node[left,font=\small] {load bias}(bp);
\draw[dashed,->] (s)--node[right,font=\small] {pointer changes}(bs);
\end{tikzpicture}
\caption{A common load bias cancels from internal PC-relative differences but remains in absolute pointer values.}
\label{fig:ch59-bias}
\end{figure}

This argument assumes that the selected symbol belongs at the assumed relative position.
If runtime symbol lookup interposes a definition in another shared object, the bases need not cancel.
Position independence and nonpreemptability are therefore separate properties.
An instruction sequence valid for a hidden local definition may be invalid for a preemptable global name even when both are located in the same input source file before linking.

\section{The dynamic section is a runtime directory}
\label{sec:ch59-dynamic-section}

The program interpreter named by \code{PT\_INTERP}, when present, participates in establishing the process image.
Dynamic objects expose tables through their dynamic metadata: dependencies, dynamic symbols and strings, hashes, relocations, and initialization information.
The runtime linker follows these loader-facing structures rather than relying on the section-header table that a static analysis tool might display.
An executable can be stripped of many section names while retaining the dynamic information needed to run.

The exact dependency search path is operating-system and runtime-linker policy.
Environment variables, embedded search paths, cache files, and privileged-execution rules are not universal ELF formulas.
A book example should name the runtime being observed instead of presenting one loader's path order as a property of all linkers.
Our executable experiment uses the Linux/POSIX \code{dlopen} interface with an explicit file path.

Table~\ref{tab:ch59-records} separates information used for static linking from information required at runtime.
The same conceptual symbol can appear in both a full symbol table and a smaller dynamic symbol table with different index spaces.
A dynamic relocation must refer to the table specified by its runtime contract, not an arbitrary index retained from a relocatable input.

\begin{table}[htbp]
\caption{Dynamic loading consumes runtime metadata rather than replaying the original static link command.}
\label{tab:ch59-records}
\centering
\begin{tabular}{p{0.29\linewidth}p{0.58\linewidth}}
\toprule
Structure & Principal role \\
\midrule
Input section table & Describes relocatable contributions for link editing. \\
Program headers & Describe mappings and runtime metadata locations. \\
Dynamic symbol table & Supplies symbols participating in runtime lookup. \\
Dynamic relocations & Describe fields still requiring runtime work. \\
Dependency entries & Name shared-object dependencies under loader policy. \\
Initialization records & Identify runtime initialization/finalization work. \\
\bottomrule
\end{tabular}
\end{table}

\section{GOT entries centralize mutable addresses}
\label{sec:ch59-got}

A global offset table can hold runtime-resolved addresses outside otherwise shareable text.
An instruction reaches a nearby table slot through a position-independent addressing sequence and loads the final pointer from that slot.
This separates the code bytes from a process-specific binding decision.
It does not eliminate relocation work: it moves that work to a data structure whose mutability and lifetime can be managed explicitly.

Consider a table slot at linked address \code{0x3000}, with a runtime definition selected at \code{0x7F0012345000}.
The runtime relocation writes that selected address into the slot.
Every use site that accesses the slot then observes the same binding.
If a call instruction instead encoded a fixed direct displacement to a different object, moving either object or changing the selected definition could invalidate the call.

The exact GOT-related relocation formulas are architecture-specific.
Some refer to the address of a symbol's slot, some to the GOT base, and some express the slot's displacement from the instruction.
The names cannot be interpreted by a generic rule such as \enquote{write the symbol address minus the table base}.
The AMD64 supplement distinguishes these values explicitly~\cite{matz-2012-amd64}.

\section{PLT binding is a state transition}
\label{sec:ch59-plt}

A procedure linkage table provides call stubs that cooperate with runtime binding.
Under a lazy-binding policy, an unresolved slot initially routes a call into the runtime resolver.
After lookup, the slot can be updated so later calls reach the selected function through a shorter path.
Under eager binding, the runtime resolves the relevant slots before transferring control to ordinary program code.
These policies move work between startup and first use; neither guarantees better total performance for every workload.

Figure~\ref{fig:ch59-lazy} depicts the state transition without committing to one processor's exact stub instruction sequence.
Concurrent first calls require synchronization consistent with the runtime's rules.
Auditing, profiling, or interposition facilities can also affect whether a slot is permanently rewritten.
An implementation should not assume that a conceptual lazy-binding diagram is a complete description of every production loader.

\begin{figure}[htbp]
\centering
\begin{tikzpicture}[node distance=16mm and 22mm]
\node[basic block] (call) {call site};
\node[basic block,right=of call] (slot) {PLT / address slot};
\node[basic block,below left=of slot] (resolver) {runtime resolver};
\node[basic block,below right=of slot] (fun) {chosen function};
\draw[cfg edge] (call)--(slot);
\draw[dashed,->] (slot)--node[left,font=\small,fill=white,inner sep=2pt] {unresolved}(resolver);
\draw[cfg edge] (resolver)--node[below,font=\small] {select and publish}(fun);
\draw[cfg edge] (slot)--node[right,font=\small,fill=white,inner sep=2pt] {resolved}(fun);
\end{tikzpicture}
\caption{Lazy binding changes a runtime indirection after lookup; eager binding performs that work before first use.}
\label{fig:ch59-lazy}
\end{figure}

\section{Runtime lookup is not static strong-over-weak merging}
\label{sec:ch59-lookup}

The static resolution table in Chapter~\ref{ch:link-symbols} gives strong definitions precedence over weak ones.
Runtime lookup additionally depends on an ordered search scope, visibility, version requirements, and the loader's binding policy.
It is not generally implemented by collecting every loaded definition and selecting the globally strongest one.
An earlier eligible runtime definition can determine a binding even when another loaded object has a different definition with the same spelling.
Platform rules must be consulted before treating static weak-symbol intuition as a dynamic lookup algorithm.

\begin{algorithm}[htbp]
\caption{Resolve a runtime reference under an explicitly supplied lookup policy.}
\label{alg:ch59-lookup}
\begin{pseudocode}
\Require name/version request, ordered lookup scope, platform eligibility and binding policy
\Ensure a permitted runtime definition or a policy-defined unresolved outcome
\ForAll{objects in the lookup scope's declared order}
  \State query the object's dynamic symbol index for the requested name
  \ForAll{candidate entries}
    \State check definition state, visibility, version, and permitted symbol kind
    \If{the policy selects this eligible candidate}
      \State \Return its runtime definition
    \EndIf
  \EndFor
\EndFor
\State apply the platform's missing-symbol rule for this reference
\end{pseudocode}
\end{algorithm}

Algorithm~\ref{alg:ch59-lookup} intentionally exposes policy as input rather than inventing a universal ELF search order.
With $D$ objects in the scope and $H$ the total work of their hash/index queries and candidate checks, its search cost is $O(D+H)$.
Real loaders cache results, use version-aware hash tables, and special-case local or relative relocations that need no global lookup.
Those fast paths preserve the same distinction between locating a name and computing a field value.

\section{Relative relocations and initialization order}
\label{sec:ch59-relocation}

A relative relocation often needs only the object's load bias and an addend.
On AMD64, \code{R\_X86\_64\_RELATIVE} computes $B+A$.
It does not require a name search for an ordinary external symbol.
Grouping or compactly representing such relocations can reduce loader overhead, but the encoding and permitted ordering are ABI details.
Table~\ref{tab:ch59-runtime-reloc} contrasts common runtime tasks.

\begin{table}[htbp]
\caption{Runtime relocations differ in whether they need symbol lookup or only a load bias.}
\label{tab:ch59-runtime-reloc}
\centering
\begin{tabular}{lll}
\toprule
AMD64 operation & Representative value & Needed information \\
\midrule
RELATIVE & $B+A$ & object's load bias \\
GLOB\_DAT & $S$ & selected symbol address \\
JUMP\_SLOT & $S$ & selected callable definition \\
TLS-related relocation & ABI-specific & module/thread-relative information \\
\bottomrule
\end{tabular}
\end{table}

After mappings and required relocation work are ready, initialization functions can run under the runtime's dependency and cycle rules.
Dependencies with cycles cannot all satisfy a strict \enquote{every dependency initializes first} ordering.
The loader needs an explicit cycle policy and must avoid reentering the same initialization indefinitely.
Finalization and unloading similarly involve dependency, reference, and lifetime rules rather than a universal reversal of the static link command.

Code that obtains a function pointer from a dynamically loaded object must not continue calling it after that object's code is unloaded.
Reference counts, retained handles, and application synchronization are therefore part of the runtime interface.
Relocation correctness alone does not establish lifetime correctness.
Chapter~\ref{ch:extension-models} turns that requirement into a versioned plugin-table contract: the retained loader handle owns the validity interval of its function entries and object destructors.
Its native test rejects incompatible table requests and unloads only after the last table use.

\section{TLS adds a thread coordinate}
\label{sec:ch59-tls}

An ordinary global object has one address within a process mapping.
A \term{thread-local storage} object has a distinct instance for each thread.
The linker can describe the variable's offset within a TLS template, but the eventual address also depends on which thread is accessing it.
The runtime must copy initialized TLS bytes and zero the remaining TLS extent for the relevant thread storage.
ELF uses a TLS program-header description for the initialization template~\cite{xinuos-2026-elf}.

Figure~\ref{fig:ch59-tls} uses a template offset of $24$.
Thread one and thread two have different TLS block bases, so the same offset identifies different addresses.
Treating the symbol as an ordinary process-wide absolute address would make the threads share storage incorrectly.
Different TLS access models trade flexibility for cheaper address formation, under assumptions about when a module becomes known and how storage is allocated.

\begin{figure}[htbp]
\centering
\begin{tikzpicture}[x=1cm,y=1cm]
\node[basic block] (template) at (0,0) {TLS variable: template +24};
\node[basic block] (a) at (-3,-1.7) {thread 1 base: 0x710000};
\node[basic block] (b) at (3,-1.7) {thread 2 base: 0x720000};
\node[basic block] (av) at (-3,-3.4) {instance: 0x710018};
\node[basic block] (bv) at (3,-3.4) {instance: 0x720018};
\draw[cfg edge] (template)--(a); \draw[cfg edge] (template)--(b);
\draw[cfg edge] (a)--(av); \draw[cfg edge] (b)--(bv);
\end{tikzpicture}
\caption{One TLS symbol offset denotes different concrete storage in different threads.}
\label{fig:ch59-tls}
\end{figure}

The runnable experiment in \path{dynamic-lib.c} defines a thread-local counter initialized to $37$ and exports a function returning and incrementing it.
The loader driver calls once in the main thread, twice in a worker thread, then once more in the main thread after joining the worker.
The observed per-thread sequences are both $37,38$.
The driver uses eager local loading, checks the dynamic lookup error state, and closes the library only after the worker has finished.
Both GCC and Clang builds execute this experiment; it uses the host runtime loader, not our static case-study linker.

\section{Link-time garbage collection follows references from roots}
\label{sec:ch59-gc}

After archive extraction, a linker can discard unreachable allocated sections under an appropriate root policy.
Relocations induce edges from a referencing section to the section containing the selected target.
The entry section, explicitly retained sections, and other ABI- or export-required sections supply roots.
Initialization arrays and externally discoverable entry points may need explicit retention even when ordinary code references do not reveal them.
GNU scripts provide \code{KEEP} for such cases~\cite{gnu-2026-ld}.

\begin{algorithm}[htbp]
\caption{Compute live input sections from a declared root and reference policy.}
\label{alg:ch59-section-gc}
\begin{pseudocode}
\Require selected objects, resolved relocation targets, complete root set and implicit edges
\Ensure every section reachable under that policy remains live
\State mark roots and place them on a worklist
\While{the worklist is nonempty}
  \State remove section $s$
  \ForAll{reference edges $s\to t$}
    \If{$t$ is unmarked}
      \State mark $t$ and add it to the worklist
    \EndIf
  \EndFor
\EndWhile
\State discard only unmarked discardable sections; recompute layout afterward
\end{pseudocode}
\end{algorithm}

With $V$ sections and $E$ reference edges, Algorithm~\ref{alg:ch59-section-gc} costs $O(V+E)$ using adjacency lists.
The difficult correctness premise is that the root/edge policy captures all required uses.
An inline assembly convention, a runtime registration table, or a symbol found by name can violate a naive \enquote{only direct calls matter} graph.
Graph traversal is easy once the graph is sound; discovering the right graph is the toolchain-specific part.

\section{Folding and LTO require stronger equivalence arguments}
\label{sec:ch59-transforms}

Identical-code folding cannot rely solely on equal unrelocated bytes.
Two byte-identical functions can contain relocation placeholders naming different callees.
Their relocation targets, addends, associated metadata, and address-significance policy can distinguish them.
Conversely, different input symbol indices can still name equivalent selected targets, so raw relocation-record bytes can be too strict a comparison.
An iterative structural comparison may be needed when candidate functions refer to other candidates.

Even code with identical behavior can have observably distinct addresses under a language or toolchain contract.
Folding those addresses requires either a permitted optimization policy or evidence that the addresses are not significant.
Debugging, unwind, exception, and profiling records must also be updated consistently.
This is another example of a transformation preserving more than just instruction bytes.

LTO moves some compiler work into the link pipeline by retaining an intermediate representation until more whole-program information is available.
It can enable cross-module inlining and dead-code elimination before machine code is finalized.
Assembly objects, opaque libraries, dynamic interposition, and externally visible symbols still constrain what that \enquote{whole program} means.
The linker must coordinate final code generation with ordinary object selection and relocation rather than treating IR files as already encoded sections.

\section{Exercises}
\label{sec:ch59-exercises}
\begin{enumerate}
\item\label{ex:ch59-bias} Add bias \code{0x400000} to place \code{0x1200} and target \code{0x1280}. Compute the old and new relative differences and the new absolute target.
\item\label{ex:ch59-tls} Compute both TLS addresses in Figure~\ref{fig:ch59-tls} for offset 40 instead of 24. Explain why one ordinary absolute relocation cannot represent both.
\item\label{ex:ch59-gc} For edges entry to a, a to b, b to a, and orphan to b, compute the live set with entry as the only root.
\item\label{ex:ch59-fold} Two functions have identical bytes with a zero call field, but their relocations name different non-equivalent targets. Explain the failure of byte-only folding.
\item\label{ex:ch59-lifetime} In the TLS experiment, explain why unloading before joining the worker can invalidate a function pointer even though dynamic symbol lookup originally succeeded.
\end{enumerate}

\section{Further reading}
\label{sec:ch59-reading}
The ELF specification describes dynamic metadata, initialization, and TLS templates~\cite{xinuos-2026-elf}.
The AMD64 supplement defines processor-specific runtime relocation and TLS mechanisms~\cite{matz-2012-amd64}.
Levine discusses dynamic linking and advanced linker techniques, while GNU documentation specifies concrete retention and script controls~\cite{levine-1999-linkers,gnu-2026-ld}.
